\chapter{Installation, binaries, and runtime topologies}\label{ch:installation}

\section{Build from source}

The release source package is a Rust workspace. The normal release build that produces the two installable programs is:

\begin{lstlisting}[language=bash]
cargo build --locked --release
\end{lstlisting}

The repository's own full self-check is \code{./verify.sh}; run it when you validate or release the RWLang source tree itself. It is not the first step required merely to use an application.

\code{--locked} ensures that the build uses the dependency graph captured in the release's \code{Cargo.lock}. If you only need the two executable programs:

\begin{lstlisting}[language=bash]
cargo build --locked --release -p rwlang-server -p rwlang-cli
\end{lstlisting}

After the build, the two important programs are:

\begin{lstlisting}
target/release/rwlang-server
target/release/rwlang-cli
\end{lstlisting}

\begin{itemize}
  \item \code{rwlang-server}: the long-running service process that compiles RWLang applications and serves HTTP(S);
  \item \code{rwlang-cli}: an operator/administrative program for tasks such as authentication administration and migrations.
\end{itemize}

Their roles are intentionally different. A reverse proxy connects only to the \code{rwlang-server} process. \code{rwlang-cli} is not a web server and does not need a public or proxied listener.

During development, these may be sufficient:

\begin{lstlisting}[language=bash]
cargo check --workspace
cargo test --workspace
\end{lstlisting}

\section{Where to install the binaries}

On a production host, do not run the service permanently out of Cargo's \code{target/} directory. RWLang is locally installed software rather than a program managed by the operating system's package manager, so this book consistently uses the \code{/usr/local} hierarchy: executables go under \code{/usr/local/bin}, and machine-local configuration goes under \code{/usr/local/etc/rwlang}. This keeps the RWLang installation separate from distribution-managed \code{/usr/bin} and \code{/etc} files.

A simple, auditable layout is:

\begin{lstlisting}
/usr/local/bin/rwlang-server
/usr/local/bin/rwlang-cli
/usr/local/etc/rwlang/server.toml
/srv/rwlang/current/
/srv/rwlang/data/
/run/secrets/rwlang/
/var/log/rwlang/
\end{lstlisting}

For example:

\begin{lstlisting}[language=bash]
sudo install -d -m 0755 /usr/local/bin
sudo install -d -m 0755 /usr/local/etc/rwlang
sudo install -m 0755 target/release/rwlang-server /usr/local/bin/rwlang-server
sudo install -m 0755 target/release/rwlang-cli    /usr/local/bin/rwlang-cli
sudo install -m 0644 config/server.toml.sample \
  /usr/local/etc/rwlang/server.toml
\end{lstlisting}

\code{/usr/local/bin} is normally already part of the system PATH, so no additional symlink is necessary. The examples in this book use the actual binary names, \code{rwlang-server} and \code{rwlang-cli}.


\subsection*{Debian package installation is intentionally different}

The layout above is for manual/source installation. If RWLang is installed through a Debian package, the files are owned by \code{dpkg}, so Debian filesystem conventions apply instead of the local-administrator \code{/usr/local} prefix:

\begin{lstlisting}
/usr/bin/rwlang-server
/usr/bin/rwlang-cli
/etc/rwlang/server.toml
/usr/lib/systemd/system/rwlang.service
\end{lstlisting}

Build the package on a Debian/Ubuntu build host with:

\begin{lstlisting}[language=bash]
make deb
sudo dpkg -i dist/rwlang_1.0.0-1_amd64.deb
\end{lstlisting}

The package deliberately does not start the service automatically: deploy the application and secrets, edit \code{/etc/rwlang/server.toml}, validate it, and only then enable \code{rwlang.service}. This is not a contradiction with the manual layout; \code{/usr/local} is correct for administrator-owned local installs, while \code{/usr} and \code{/etc} are correct for distribution/package-manager-owned files. The repository's \code{docs/36-debian-package.md} gives the complete packaging workflow.

\section{Your first application}

A minimal recommended project is:

\begin{lstlisting}
myapp/
|-- main.rw
|-- pages.rw
`-- public/
\end{lstlisting}

\code{main.rw}:

\begin{lstlisting}[language=RWLang]
mod pages;
\end{lstlisting}

\code{pages.rw}:

\begin{lstlisting}[language=RWLang]
page fn home(ctx: PageContext) -> Result<Html, PageError> {
    return Ok(html {
        <main><h1>Hello RWLang</h1></main>
    });
}

route home GET "/" => home;
\end{lstlisting}

Start it directly from the workspace during development:

\begin{lstlisting}[language=bash]
cargo run -p rwlang-server -- \
  --app /abs/path/myapp/main.rw \
  --listen 127.0.0.1:8080 \
  --insecure-dev-cookies
\end{lstlisting}

The same command using the installed release binary is:

\begin{lstlisting}[language=bash]
/usr/local/bin/rwlang-server \
  --app /abs/path/myapp/main.rw \
  --listen 127.0.0.1:8080 \
  --insecure-dev-cookies
\end{lstlisting}

\code{--insecure-dev-cookies} is for local development only.

\section{At first, learn only the minimum configuration}

At this point in the book the goal is not to learn the entire \code{server.toml}. As a beginner, it is enough to distinguish three things:

\begin{itemize}
  \item the \emph{application code} defines the models, queries, pages, and routes that exist;
  \item the \emph{server configuration} defines where the process listens, which external resources it may access, and what operator-controlled limits apply;
  \item \emph{secret files} hold sensitive values that must not live in application source or plaintext TOML.
\end{itemize}

The A-to-Z chapter provides a minimal development configuration. Production-specific keys are introduced only after there is an application worth operating; the complete key/default/sample reference is near the end of the book. On a first reading, do not try to memorize the full configuration surface.

\section{Production process and reverse proxy}

Behind Apache or Nginx, the typical request path is:

\begin{lstlisting}
Internet
   |
   | HTTPS :443
   v
Apache or Nginx
   |
   | HTTP 127.0.0.1:8080
   v
/usr/local/bin/rwlang-server
   |
   +-- /usr/local/etc/rwlang/server.toml
   +-- /srv/rwlang/current/main.rw
   +-- DB / Redis / AppFs
\end{lstlisting}

The proxy neither starts nor replaces the RWLang process. Run \code{rwlang-server} as a separate service, for example under systemd, and let the proxy forward requests to its loopback listener. \code{rwlang-cli} runs only for operator commands.

A simple systemd \code{ExecStart} fragment is:

\begin{lstlisting}
ExecStart=/usr/local/bin/rwlang-server \
  --config /usr/local/etc/rwlang/server.toml
\end{lstlisting}

The process user should not be \code{root}; grant access only to the required application, secret, data, and log files. If TLS is terminated by Apache or Nginx, keep the RWLang listener on loopback, for example \code{127.0.0.1:8080}.

\section{Runtime topologies in detail}

The binary locations and process roles stay the same whether RWLang terminates TLS itself or runs behind Apache or Nginx. The difference is the listener, TLS termination, and trusted-proxy boundary configuration.

\section{Direct single-domain HTTPS}

When RWLang terminates TLS itself, a public listener may use port 443:

\begin{lstlisting}
[server]
app = "/srv/rwlang/current/main.rw"
listen = "0.0.0.0:443"
insecure_dev_cookies = false

[tls]
cert_file = "/run/secrets/rwlang/tls-fullchain.pem"
key_file = "/run/secrets/rwlang/tls-key.pem"
handshake_timeout_ms = 5000
http_redirect_listen = "0.0.0.0:80"
public_host = "www.example.com"
\end{lstlisting}

\code{public_host} comes from trusted configuration. Only GET/HEAD requests may be redirected from HTTP to HTTPS.

Because this direct configuration binds ports 80/443, an unprivileged systemd service needs the narrow \code{CAP_NET_BIND_SERVICE} capability. Do not run the service as root merely to bind these ports. Behind a reverse proxy, with a listener such as \code{127.0.0.1:8080}, that capability is unnecessary.

\section{Behind an Apache reverse proxy}

In this topology Apache handles public TLS and RWLang listens only on loopback. Example RWLang configuration:

\begin{lstlisting}
[server]
app = "/srv/rwlang/current/main.rw"
listen = "127.0.0.1:8080"
insecure_dev_cookies = false

[tls]
public_host = "www.example.com"

[web]
trusted_proxy_cidrs = ["127.0.0.1/32", "::1/128"]
allow_missing_origin = false
cors_origins = []
cors_allow_credentials = false
\end{lstlisting}

In reverse-proxy mode there is no RWLang-side certificate/key, but \code{tls.public_host} is still required. It is the trusted source for the canonical host visible to the browser and anchors Host and Origin/Referer validation. Production plain-HTTP upstream is allowed only on a loopback listener with an explicit trusted-proxy list and such a public host. Every normal proxied request must carry trusted metadata such as \code{X-Forwarded-Proto: https} or the equivalent standard \code{Forwarded: proto=https} from a trusted proxy.

A minimal Apache VirtualHost is:

\begin{lstlisting}
<VirtualHost *:443>
    ServerName www.example.com

    SSLEngine on
    SSLCertificateFile /etc/letsencrypt/live/www.example.com/fullchain.pem
    SSLCertificateKeyFile /etc/letsencrypt/live/www.example.com/privkey.pem

    ProxyPreserveHost On
    ProxyPass        / http://127.0.0.1:8080/
    ProxyPassReverse / http://127.0.0.1:8080/

    RequestHeader set X-Forwarded-Proto "https"
</VirtualHost>
\end{lstlisting}

This typically requires the \code{proxy}, \code{proxy_http}, \code{headers}, and \code{ssl} modules. RWLang accepts forwarding headers only from ranges listed in \code{trusted_proxy_cidrs}; do not configure an overly broad range such as the entire Internet.

The port-80 VirtualHost may exist solely to redirect to HTTPS:

\begin{lstlisting}
<VirtualHost *:80>
    ServerName www.example.com
    Redirect permanent / https://www.example.com/
</VirtualHost>
\end{lstlisting}

\section{Behind an Nginx reverse proxy}

The Nginx topology uses the same trust model: Nginx handles public TLS and RWLang listens on loopback. The RWLang-side \code{server.toml} may be the same as in the Apache example.

A minimal Nginx server block is:

\begin{lstlisting}
server {
    listen 443 ssl;
    server_name www.example.com;

    ssl_certificate     /etc/letsencrypt/live/www.example.com/fullchain.pem;
    ssl_certificate_key /etc/letsencrypt/live/www.example.com/privkey.pem;

    location / {
        proxy_pass http://127.0.0.1:8080;
        proxy_http_version 1.1;
        proxy_set_header Host $host;
        proxy_set_header X-Forwarded-Proto https;
        proxy_set_header X-Forwarded-For $remote_addr;
    }
}
\end{lstlisting}

This example has a single direct reverse proxy, so \code{X-Forwarded-For} is the client address as seen by the proxy and \code{X-Forwarded-Proto} is \code{https}. Do not let arbitrary forwarding headers arriving from the Internet become authority. RWLang only trusts forwarding headers received from the exact \code{trusted_proxy_cidrs} ranges.

The HTTP listener may exist only to redirect to HTTPS:

\begin{lstlisting}
server {
    listen 80;
    server_name www.example.com;
    return 301 https://$host$request_uri;
}
\end{lstlisting}

\section{Recommended production directory structure}

\begin{lstlisting}
/usr/local/bin/rwlang-server
/usr/local/bin/rwlang-cli
/srv/rwlang/releases/2026-09-04.1/
/srv/rwlang/current -> /srv/rwlang/releases/2026-09-04.1
/srv/rwlang/data/
/usr/local/etc/rwlang/server.toml
/run/secrets/rwlang/
/var/log/rwlang/
\end{lstlisting}

The binary, release tree, and configuration should be runtime read-only. The application should write only to declared data/log locations.

\section{First validation}

Before sending traffic to the server:

\begin{lstlisting}[language=bash]
/usr/local/bin/rwlang-server \
  --config /usr/local/etc/rwlang/server.toml \
  --check-config
\end{lstlisting}

This validates configuration, secret-file references, application compilation, and TLS files, among other checks, without opening a listener.
